\section{Análise dos resultados}

As rampas apresentam amplitudes harmônicas proporcionais a $1/n$,
enquanto os triângulos têm decaimento proporcional a $1/n^2$ e
coeficientes pares nulos. Assim, a mesma ordem máxima 7 aproxima
melhor as ondas triangulares, já que os valores tendem a ser menores. No triângulo bipolar, a fundamental
contém aproximadamente $98{,}5534\%$ da energia; com as ordens até 7,
essa fração sobe para $99{,}9689\%$.

Nas rampas e no triângulo unipolar, a média $A/2$ representa
$75\%$ da energia total. Dessa forma, a rampa unipolar tem apenas
$22{,}9766\%$ nas ordens 1 a 7, embora a reconstrução que inclui
a média retenha $97{,}9766\%$.

Além disso, mover uma retangular dentro de um período, mantendo $D$, altera
$a_n$, $b_n$ e a fase, mas não $C_n$ nem as energias harmônicas.
Os casos simétrico e assimétrico confirmam essa propriedade.
A retangular bipolar só possui média zero para $D=1/2$; nos exemplos
com $D=1/3$, sua média é $-A/3$.

As seis ondas deslocadas obedecem à mesma invariância de energia.
Já alterar $\delta$ modifica a posição das linhas em hertz:
os exemplos de 60 e 75 Hz precisam de janelas de duração $1/60$ e
$1/75$ s, respectivamente. Usar o período nominal $P$ para analisar
uma onda com período efetivo diferente pode causar vazamento
espectral e comparar frequências que não representam as mesmas ordens.

Na alternância dos dois pulsos, o intervalo repetido é $2P$ e a
fundamental do exemplo passa a 30 Hz. A fração de energia nessa
ordem cai para $6{,}6934\%$ no unipolar e $11{,}1557\%$ no bipolar.

As reconstruções com saltos apresentam oscilações próximas às
descontinuidades, associadas ao fenômeno de Gibbs. Esse erro de
truncamento é diferente do pequeno erro da FFT amostrada em relação
aos coeficientes analíticos.




